\section{Subquadratic interior saddle with a fifth-order remainder}\label{sec:subquadratic-interior-saddle}

Let $d\to\infty$, $1\le r<d$, and let $S$ be a binomial law tilted and truncated to $\{0,\ldots,r\}$, with weights proportional to $\binom ds v^s$, $v>0$. Set $\mu=\mathbb ES$, $V=\mathrm{Var}(S)$, $a_s=s(d-s)/[d(d-1)]$, and $a=\mathbb E a_S$. For a uniform $S$-subset $I$, use $R(\theta)=\mathbb E\exp(i\sum_{i\in I}\theta_i)$ and $\phi(t)=\mathbb E e^{itS}$. Write $H=\sum_{s\le r}\binom ds v^s$, $a(t)=\mathbb E[a_S e^{itS}]/\phi(t)$, and $t_*=\pi/d$. Fix $h_0>0$ and assume

\begin{equation*}
\mu\ge1/10,\quad r-\mu\ge h_0,\quad
n/d^{5/3}\to\infty,\quad na/d\to\infty,\quad nV/d^2\to\infty.
\end{equation*}

Then $P_n^{(r,d)}(-v^d)/H^n$ equals

\begin{equation*}
\frac{2}{\sqrt d}(2\pi n)^{-(d-1)/2}\mathrm{Re}\!\left[\phi(t_*)^n a(t_*)^{-(d-1)/2}\right]+o(B_*),
\qquad B_*=\frac{|\phi(t_*)|^n}{\sqrt d}(2\pi n|a(t_*)|)^{-(d-1)/2}.
\end{equation*}

The error is absolute on the scale of one principal saddle, uniformly along the stated sequences. The conclusion concerns the angular integral on this interior range.

\subsection{Pointwise fifth-moment estimate}

For $y\in\mathbf1^\perp$ put $L=\sum_{i\in I}y_i$, $Q_j=\sum_i y_i^j$. At fixed size $s$ the centered vector allows passage to the complement with a sign reversal, so use $m=\min(s,d-s)$ draws. Hoeffding's convex comparison theorem for sampling without replacement, with $\Phi(x)=|x|^5$, bounds $\mathbb E[|L|^5\mid S=s]$ by the corresponding moment for $m$ independent draws from the finite population $y_1,\ldots,y_d$. This is Theorem 4 of Hoeffding (1963). \href{https://doi.org/10.1080/01621459.1963.10500830}{Published DOI}, \href{https://www.cs.rpi.edu/academics/courses/spring06/random/hoefding.pdf}{scan of primary article}.

The comparison can also be seen by drawing an ordered sample of $m$ distinct population indices, independently drawing the equality pattern of $m$ draws with replacement, and assigning the pattern blocks in first-appearance order to the first distinct sampled indices. The resulting sequence has the with-replacement law. Conditional on the unordered distinct sample and the pattern, each assigned population value has the sample average as its expectation. The conditional expected sum is therefore the without-replacement sample sum. Jensen gives the desired convex comparison.

Here is a short proof of the independent fifth-moment bound, avoiding a separate moment-theorem dependency. For centered independent identically distributed draws $X_i$ with variance $\sigma_0^2$, $S_j'=\sum_{i\le j}X_i$, and $M_j=\mathbb E|S_j'|^5$, Taylor gives
$|x+u|^5\le |x|^5+5x|x|^3u+C(|x|^3u^2+|u|^5)$. The linear term vanishes upon conditioning on the past. Jensen gives $M_{j-1}\le M_j$. Summing the moment increments yields $M_m\le C m\sigma_0^2 M_m^{3/5}+Cm\mathbb E|X_1|^5$. Splitting according to which right-hand term is larger implies $M_m\le C[(m\sigma_0^2)^{5/2}+m\mathbb E|X_1|^5]$. Here $\sigma_0^2=Q_2/d$ and $\mathbb E|X_1|^5=d^{-1}\sum_i|y_i|^5$. Since $m/d\asymp a_s$ for $0<s<d$, we obtain

\begin{equation*}
\mathbb E[|L|^5\mid S=s]\le C\left(a_s\sum_i|y_i|^5+a_s^{5/2}Q_2^{5/2}\right).
\end{equation*}

Both sides vanish for $s=0$. The moment bound $\mathbb E a_S^{5/2}\le C a^{5/2}$ follows from $O(1/d)$ Lipschitz continuity, $a_\mu=a+V/[d(d-1)]$, $a\ge\max(V/d,1/(11d))$, and $\mathbb E|S-\mu|^{5/2}\le C(1+V)^{5/4}$. The latter follows from log-concavity of $S-S'$ and its geometric survival bound at scale $\sqrt{1+V}$, as proved in \href{https://github.com/ProgramComputer/open-math-research/blob/4fb624ba937a8eecc1361a6066185cd81ab7b172/topics/balanced-arrays/quadratic/very-tall-quadratic-interior-saddle.md}{the quadratic interior note}. Thus the unconditional fifth bound has the same form with $a$ replacing $a_s$.

\subsection{Near-chart expansion and exponential domination}

Use the angular charts $\theta=t\mathbf1+y$ with $\|y\|\le\eta$ at odd multiples $t$ of $\pi/d$. For $|t|\le c_0/\sqrt{1+V}$, $\phi(t)$ is bounded away from zero. The normalized complex row logarithm is

\begin{equation*}
\log\!\left(\frac{R(t\mathbf1+y)}{\phi(t)}\right)
=-\tfrac12a(t)Q_2-\tfrac i6\gamma(t)Q_3+K_4(y)+B_5(y).
\end{equation*}

Here $|\gamma(t)|\le Ca$, $|\mathrm{Im}\gamma(t)|\le Ca\sqrt{1+V}|t|$, and $|K_4|\le CaQ_4$, using the cancellation in \href{https://github.com/ProgramComputer/open-math-research/blob/605cc467a5be93a8408519da331e46f1008df96e/topics/balanced-arrays/critical-aspect/fourth-transverse-cumulant.md}{the fourth cumulant proof}. For clarity, the $Q_2^2$ coefficient left after taking the logarithm is $O(a/d)$: at fixed $s$ the raw coefficient divided by 3 is $(s)_2(d-s)_2/(d)_4=a_s^2+O(a_s/d)$, and the complex variance of $a_S$ is $O(V/d^2)\le O(a/d)$. Then $Q_2^2/d\le Q_4$.

The log remainder satisfies $|B_5|\le C\mathbb E|L|^5$. Indeed $R/\phi-1=O(\mathbb E|L|^2)$ because the conditional first moment vanishes. Taylor of the exponential through order four gives a remainder bounded by the fifth absolute moment. For the logarithm, only the square of the quadratic term contributes at order four. Bounding the difference from that square uses $O((\mathbb E|L|^2)(\mathbb E|L|^3))$ and the log cubic remainder uses $O((\mathbb E|L|^2)^3)$. Both are at most $C\mathbb E|L|^5$ by moment monotonicity and the bound $\mathbb E|L|^2\le C\eta^2$. Consequently $|B_5|\le C(a\sum_i|y_i|^5+a^{5/2}Q_2^{5/2})$.

The diagonal coefficient estimates from the quadratic note give $\alpha_t=\mathrm{Re}\,a(t)\asymp a$, $\alpha_t/a=1+O((1+V)t^2)$ and $\mathrm{Im}\,a(t)/a=O(\sqrt{1+V}|t|)$. Let $A=n\alpha_t\asymp na$. The real exponential correction above the Gaussian of precision $A$ is bounded on the chart by

\begin{equation*}
Cna\left(\sqrt{1+V}|t|\sum_i|y_i|^3+Q_4\right)
+Cna^{5/2}Q_2^{5/2},
\end{equation*}

using $\sum |y_i|^5\le\eta Q_4$. Its exponential moment, with any fixed multiplier involved in Cauchy--Schwarz, is $O(\exp(Cd(1+V)t^2))$. To verify this, separate coordinate and radial terms by Cauchy--Schwarz. For the coordinate terms represent $y_i=Z_i-\bar Z$ with iid $N(0,A^{-1})$ and restrict also to $|\bar Z|\le\eta$, independently with probability $1-o(1)$. Then $|Z_i|\le2\eta$ and $\sum |y_i|^j\le2^j\sum |Z_i|^j$. On the resulting interval the single-coordinate exponent can be absorbed into half the original Gaussian quadratic by making $\eta$ small. Integrating $e^u-1\le u e^u$ yields a factor at most $1+C(\sqrt{1+V}|t|/\sqrt A+1/A)$ per coordinate. Their product is at most $\exp(Cd((1+V)t^2+1/A))$ by Young's inequality, with $d/A=o(1)$.

For the radial term put $U=AQ_2$, a chi-square variable with $d-1$ degrees of freedom. The exponent is at most $C U^{5/2}/n^{3/2}$. On $U\le L_0 d$ for fixed large $L_0$, this is $O_{L_0}(d^{5/2}/n^{3/2})=o(1)$. For $U>L_0d$ on the chart it is at most $\epsilon U$ with fixed $\epsilon<1/8$, since $Q_2\le\eta^2$ and $a\le1$. The chi-square exponential moment at $1/4$ bounds this part by $\exp(-c d)$ after taking $L_0$ sufficiently large. This proves the exponential domination without an $\exp(O(d^2/n))$ penalty.

\subsection{Integrated comparison and phase exclusion}

Under the same real Gaussian, the $L^2$ norms of $Q_3$, $Q_4$ and $\sum_i|y_i|^5$ are respectively $O(\sqrt d\,A^{-3/2})$, $O(d A^{-2})$, and $O(d A^{-5/2})$. The first is Wick's formula with the row-sum-zero covariance cancellation; the others follow by the triangle inequality. Also $\|Q_2^{5/2}\|_2=O((d/A)^{5/2})$. Multiplying by the respective coefficients and $n$ bounds the total log perturbation in $L^2$ by

\begin{equation*}
O\!\left(\sqrt{\frac d{na}}+\frac d{na}+\frac{d^{5/2}}{n^{3/2}}\right)=o(1).
\end{equation*}

The coordinate quintic term is absorbed into the quartic bound on the chart. Applying Cauchy--Schwarz to the difference of exponentials gives an absolute chart error bounded by this $o(1)$ times $C I_0|\phi(t)|^n e^{Cd(1+V)t^2}$, where $I_0=(2\pi na)^{-(d-1)/2}/\sqrt d$. The complex Gaussian may be extended beyond the chart because $na\gg d$. The chart modulus itself is at most $C I_0|\phi(t)|^n e^{Cd(1+V)t^2}$. At $t_*$ the determinant correction is bounded because $V=O(d)$, and hence the two principal chart errors are $o(B_*)$.

For completeness the global localization uses only the unchanged hypotheses behind the quadratic note: off-chart the random disjoint-pair estimate yields $|R|\le\exp[-c\min(aD_z,1)]$, $D_z=\sum_i|z_i-\bar z|^2$. On the unsaturated part, retaining half the penalty and integrating the other half using Haar relative angles at a closest anchor gives relative bound $\exp(O(d)-cna+O(nV/d^2))$. This vanishes because $na\gg d$ and $V/d^2\le a/d$. The saturated part and far diagonal charts have row modulus bounded away from one, so their suppression $e^{-cn}$ absorbs the reciprocal Gaussian factor $e^{O(d\log n)}$ and radial loss $e^{O(nV/d^2)}$. Here $n\gg d\log n$. The normalized torus measure, of Euclidean covolume $\sqrt d(2\pi)^{d-1}$, gives the factor in the Gaussian formula.

Finally the near radial gap is $|\phi(t)/\phi(t_*)|\le e^{-c V(t^2-t_*^2)}$ for $t_*\le|t|\le c_0/\sqrt{1+V}$, by the fourth moment bound for $S-S'$ and differentiation as in the quadratic note. Since $V\ge c_{h_0}$ and $n\gg d$, the determinant factor $e^{Cd(1+V)t^2}$ is absorbed in this gap at odd labels $j=|t|/t_*\ge3$. The sum over such phases is $O(\sum_{j\ge3}\exp[-cnV(j^2-1)/d^2])=o(1)$ relative to $B_*$. Only the two conjugate principal charts remain.

\subsection{Raw remainder at the further four-thirds scale}

At the further height scale $n\asymp d^{4/3}$, using a raw sixth moment after retaining the signed fifth term would still leave a radial bound of order $d^3/n^2$. For instance at tilt $v=1$ and cap $r=d-1$, the transverse coefficient $a$ is of constant order. On a typical Gaussian shell $Q_2\asymp d/n$, so $n\mathbb E|L|^6\ge n(\mathbb E L^2)^3\asymp d^3/n^2$. This is a limitation of that remainder estimate. The five-thirds proof above uses a fifth absolute moment after retaining the fourth-order logarithm.
